\chapter{Morphological Taxonomy}

\chapterepigraph{%
  Classification is not the end of understanding.\\
  It is the beginning.\\
  A good classification reveals\\
  why the things in one class\\
  behave the way they do.
}{Flyxion, \textit{Semantic Infrastructure}}

\sidenote{Connects to\\canonical forms\\(Ch.\,4) and\\constraint spaces\\(App.\,A)}

Chapter~44 examined structural recurrence — how the same structure
appears across domains. This chapter examines the formal theory of
classification itself: what makes a good taxonomy, how shapes and
functions constrain each other, and how the constraint framework
provides a principled basis for classification.

\section{Shape, Function, Behavior, Constraint}

Traditional classification schemes have emphasized shape (morphology),
function (what something does), behavior (how it acts), and increasingly,
genetic or compositional structure. Each reveals different aspects of
the underlying reality.

\begin{definition}{Morphological Equivalence Class}{morph-equiv}
  Two systems $s_1, s_2$ belong to the same \emph{morphological
  equivalence class} if there exists an isomorphism $\phi : s_1 \to s_2$
  that preserves the relevant structural features (shape, topology,
  connectivity, constraint structure). The equivalence class $[s]$ is
  the morphological type.
\end{definition}

The ``relevant structural features'' depend on the classification goal.
Anatomical classification preserves gross morphology (bone structure,
organ arrangement). Functional classification preserves input-output
behavior. Constraint classification preserves the constraint structure
— which configurations are admissible and which are not.

\section{When Categories Fail}

Every classification system fails at boundaries. The species concept
fails at ring species (populations that form a continuous chain with
fertile hybridization but whose ends cannot interbreed). The
country concept fails at disputed territories. The ``word'' concept
fails at clitics, compounds, and phrasal idioms.

\begin{definition}{Category Boundary Failure}{cat-boundary}
  A classification system $\mathcal{T}$ exhibits \emph{boundary failure}
  at item $x$ if: (1) $x$ partially satisfies the criteria for
  category $C_1$ and partially satisfies the criteria for category $C_2$,
  and (2) the criteria for $C_1$ and $C_2$ are jointly unsatisfiable
  for any single item. The failure reveals that $\mathcal{T}$'s
  categories are not constraint-closed.
\end{definition}

The resolution in each case is the same: recognize that the boundary
case reveals a hidden dimension — a feature that the binary
categorization was suppressing. Ring species reveal the temporal
dimension of speciation; disputed territories reveal the political
dimension of geography; clitics reveal the prosodic dimension of
word structure.

\section{Functional Classification}

The RSVP framework suggests a classification principle that unifies
shape, function, behavior, and constraint: classify by the accessibility
landscape. Systems with similar accessibility landscapes — similar
constraint structures, similar gradient flows, similar attractors —
are in the same morphological class regardless of their physical substrate.

\begin{namedthm}{The Functional Morphology Principle}
  The deepest classification of a system is by its accessibility
  landscape: the structure of its constraint space, admissible states,
  and gradient flow. Systems that differ in physical substrate but
  share an accessibility landscape are functionally equivalent and
  will exhibit analogous behaviors under analogous perturbations.
\end{namedthm}

\begin{exercise}
  The RSVP framework classifies physical, cognitive, and social systems
  by their accessibility landscapes. Give one example of each pair:
  (a) two physical systems with similar accessibility landscapes,
  (b) two cognitive systems with similar accessibility landscapes,
  (c) one physical and one cognitive system with similar accessibility
  landscapes. For each pair, identify a specific behavior predicted
  to be analogous.
\end{exercise}

\begin{exercise}
  Propose a classification of mathematical proof strategies
  (induction, contradiction, construction, pigeonhole, probabilistic
  method, etc.) based on their constraint-space structure.
  What is the admissible region for each strategy?
  What accessibility landscape does each navigate?
\end{exercise}
